%%
%% part5-examples/ch27-math-book.tex
%%
%% Chapter 27: A Math Book — A complete sample chapter from
%% a mathematics textbook, demonstrating all MatinBook
%% features in a real-world context.
%%

\chapter{کتاب ریاضی}
\label{chap:math-book}

\section{چرا این موضوع مهم است؟}
\label{sec:math-book-why}

در فصل‌های قبل، با تمام قابلیت‌های \lr{MatinBook}
آشنا شدید. ولی دانستن هر قابلیت به‌تنهایی کافی
نیست؛ مهم این است که بتوانید همه‌ی آن‌ها را
در یک کتاب واقعی ترکیب کنید.

در این فصل، یک فصل کامل از یک کتاب ریاضی
خیالی را ارائه می‌دهیم که همه‌ی اجزای
\lr{MatinBook} را در بافت واقعی نشان می‌دهد.
هدف، الگوبرداری برای نوشتن کتاب ریاضی خودتان
است.

\mnote{این فصل، یک «مطالعه‌ی موردی» است. به‌جای توضیح قابلیت‌ها، آن‌ها را در عمل می‌بینید.}

\section{ساختار فصل نمونه}
\label{sec:math-book-structure}

فصل نمونه، از یک کتاب ریاضی خیالی با عنوان
«مقدمه‌ای بر آنالیز حقیقی» است. این فصل
شامل بخش‌های زیر است:

\begin{itemize}
    \item مقدمه.
    \item تعاریف.
    \item قضایا و اثبات‌ها.
    \item مثال‌ها.
    \item تمرین‌ها.
    \item راه‌حل‌ها.
    \item نمودارها.
\end{itemize}

\mnote{این فصل را می‌توانید به‌عنوان الگو برای نوشتن فصل‌های کتاب ریاضی خودتان استفاده کنید.}

\section{مقدمه}
\label{sec:math-book-intro}

آنالیز حقیقی\index{آنالیز حقیقی}، شاخه‌ای از
ریاضیات است که به مطالعه‌ی اعداد حقیقی،
دنباله‌ها، سری‌ها، حد، پیوستگی و مشتق
می‌پردازد. این شاخه، پایه‌ی بسیاری از
شاخه‌های دیگر ریاضیات است.

در این فصل، با مفاهیم اساسی آنالیز حقیقی
آشنا می‌شویم و قضایای مهم آن را بررسی
می‌کنیم.

\mnote{آنالیز حقیقی، پیش‌نیاز آنالیز مختلط، آنالیز تابعی و نظریه‌ی اندازه است.}

\section{دنباله‌ها و همگرایی}
\label{sec:math-book-sequences}

\subsection{تعریف دنباله}
\label{sec:math-book-sequence-def}

\begin{definition}[دنباله]
    یک \textbf{دنباله}\index{دنباله} در
    مجموعه‌ی \lr{$\R$}، تابعی است از
    \lr{$\N$} به \lr{$\R$}:
    \[
        a: \N \to \R, \quad n \mapsto a_n
    \]
    که در آن \lr{$a_n$} را جمله‌ی \lr{$n$}-ام
    دنباله می‌نامیم.
    \label{def:sequence}
\end{definition}

طبق \cref{def:sequence}، یک دنباله را
می‌توان به‌صورت \lr{$(a_n)_{n=1}^{\infty}$}
نمایش داد.

\begin{example}
    دنباله‌ی \lr{$a_n = \frac{1}{n}$}:
    \[
        1, \frac{1}{2}, \frac{1}{3}, \frac{1}{4}, \ldots
    \]
    این دنباله به صفر نزدیک می‌شود، ولی
    هرگز به آن نمی‌رسد.
    \label{ex:harmonic}
\end{example}

\subsection{همگرایی}
\label{sec:math-book-convergence}

\begin{definition}[همگرایی دنباله]
    دنباله‌ی \lr{$(a_n)$} به عدد \lr{$L$}
    \textbf{همگرا}\index{همگرایی} است اگر:
    \[
        \forall \varepsilon > 0, \;
        \exists N \in \N, \;
        \forall n \geq N: \abs{a_n - L} < \varepsilon
    \]
    در این صورت می‌نویسیم
    \lr{$\lim_{n \to \infty} a_n = L$}.
    \label{def:convergence}
\end{definition}

\mnote{تعریف همگرایی، یکی از مهم‌ترین تعاریف در آنالیز حقیقی است. این تعریف، مفهوم «نزدیک شدن» را دقیق می‌کند.}

\section{قضایای مهم}
\label{sec:math-book-theorems}

\subsection{یکتایی حد}
\label{sec:math-book-unique-limit}

\begin{theorem}[یکتایی حد]
    اگر دنباله‌ای همگرا باشد، حد آن یکتاست.
    \label{thm:unique-limit}
\end{theorem}

\begin{proof}
    فرض کنید \lr{$\lim a_n = L_1$} و
    \lr{$\lim a_n = L_2$} با
    \lr{$L_1 \neq L_2$}. قرار دهید
    \lr{$\varepsilon = \frac{\abs{L_1 - L_2}}{2} > 0$}.

    طبق \cref{def:convergence}، عدد \lr{$N_1$}
    وجود دارد که برای \lr{$n \geq N_1$}،
    \lr{$\abs{a_n - L_1} < \varepsilon$}.
    همچنین عدد \lr{$N_2$} وجود دارد که برای
    \lr{$n \geq N_2$}،
    \lr{$\abs{a_n - L_2} < \varepsilon$}.

    برای \lr{$n \geq \max(N_1, N_2)$}:
    \begin{align*}
        \abs{L_1 - L_2}
        &= \abs{(L_1 - a_n) + (a_n - L_2)} \\
        &\leq \abs{L_1 - a_n} + \abs{a_n - L_2} \\
        &< \varepsilon + \varepsilon = \abs{L_1 - L_2}
    \end{align*}
    که تناقض است. بنابراین \lr{$L_1 = L_2$}.
\end{proof}

\subsection{قضیه‌ی فشردگی}
\label{sec:math-book-squeeze}

\begin{theorem}[قضیه‌ی فشردگی]
    اگر \lr{$a_n \leq b_n \leq c_n$} برای
    همه‌ی \lr{$n \geq N_0$} و
    \lr{$\lim a_n = \lim c_n = L$}، آنگاه
    \lr{$\lim b_n = L$}.
    \label{thm:squeeze}
\end{theorem}

\begin{proof}
    طبق \cref{def:convergence}، برای هر
    \lr{$\varepsilon > 0$}، اعداد \lr{$N_1$}
    و \lr{$N_2$} وجود دارند که:
    \begin{align*}
        n \geq N_1 &\implies \abs{a_n - L} < \varepsilon \\
        n \geq N_2 &\implies \abs{c_n - L} < \varepsilon
    \end{align*}

    برای \lr{$n \geq \max(N_0, N_1, N_2)$}:
    \[
        L - \varepsilon < a_n \leq b_n \leq c_n < L + \varepsilon
    \]
    بنابراین \lr{$\abs{b_n - L} < \varepsilon$}،
    که نتیجه می‌دهد \lr{$\lim b_n = L$}.
\end{proof}

\begin{remark}
    قضیه‌ی فشردگی (\cref{thm:squeeze}) یکی از
    ابزارهای قدرتمند برای محاسبه‌ی حد
    دنباله‌های پیچیده است. برای مثال، برای
    محاسبه‌ی \lr{$\lim \frac{\sin n}{n}$}،
    از نامساوی \lr{$-\frac{1}{n} \leq \frac{\sin n}{n} \leq \frac{1}{n}$}
    استفاده می‌کنیم.
    \label{rem:squeeze-application}
\end{remark}

\subsection{قضیه‌ی بولتزانو-وایرشتراس}
\label{sec:math-book-bolzano}

\begin{theorem}[بولتزانو-وایرشتراس]
    هر دنباله‌ی کراندار در \lr{$\R$} دارای
    زیردنباله‌ی همگراست.
    \label{thm:bolzano-weierstrass}
\end{theorem}

\mnote{قضیه‌ی بولتزانو-وایرشتراس، یکی از نتایج بنیادین در آنالیز حقیقی است. اثبات آن از اصل کامل بودن \lr{$\R$} استفاده می‌کند.}

\section{نمودارها}
\label{sec:math-book-figures}

نمودار \cref{fig:convergence-plot} همگرایی
دنباله‌ی \lr{$a_n = \frac{1}{n}$} را نشان
می‌دهد.

\begin{figure}[h]
\centering
\begin{latin}
\begin{tikzpicture}
  \begin{axis}[
    width=0.8\textwidth, height=6cm,
    xlabel={$n$},
    ylabel={$a_n$},
    grid=major,
    domain=1:20,
    samples=20,
    ymin=0, ymax=1.1
  ]
    \addplot[only marks, mark=*, color=matinblue, mark size=2pt]
        {1/x};
    \addlegendentry{$a_n = 1/n$}

    \addplot[matinred, thick, dashed, domain=1:20]
        {0};
    \addlegendentry{$\lim a_n = 0$}
  \end{axis}
\end{tikzpicture}
\end{latin}
\caption{همگرایی دنباله‌ی \lr{$a_n = \frac{1}{n}$} به صفر}
\label{fig:convergence-plot}
\end{figure}

\mnote{نمودار \lr{\cref{fig:convergence-plot}} به‌وضوح نشان می‌دهد که دنباله به صفر نزدیک می‌شود ولی هرگز به آن نمی‌رسد.}

\section{تمرین‌ها}
\label{sec:math-book-exercises}

\begin{exercise}
    ثابت کنید دنباله‌ی \lr{$a_n = \frac{n}{n+1}$}
    به \lr{$1$} همگراست.
    \label{exc:convergence-n-over-n1}
\end{exercise}

\begin{solution}
    باید نشان دهیم که برای هر
    \lr{$\varepsilon > 0$}، عدد \lr{$N$}
    وجود دارد که برای \lr{$n \geq N$}:
    \[
        \abs{\frac{n}{n+1} - 1} = \frac{1}{n+1} < \varepsilon
    \]
    با انتخاب \lr{$N > \frac{1}{\varepsilon} - 1$}،
    این نامساوی برقرار است.
\end{solution}

\begin{exercise}
    با استفاده از \cref{thm:squeeze}، حد
    دنباله‌ی \lr{$a_n = \frac{\cos n}{n}$} را
    محاسبه کنید.
    \label{exc:squeeze-cos}
\end{exercise}

\begin{solution}
    از نامساوی
    \lr{$-\frac{1}{n} \leq \frac{\cos n}{n} \leq \frac{1}{n}$}
    و \cref{thm:squeeze} استفاده می‌کنیم.
    چون \lr{$\lim \frac{1}{n} = \lim \left(-\frac{1}{n}\right) = 0$}،
    داریم \lr{$\lim \frac{\cos n}{n} = 0$}.
\end{solution}

\begin{exercise}
    آیا دنباله‌ی \lr{$a_n = (-1)^n$} همگراست؟
    چرا؟
    \label{exc:alternating}
\end{exercise}

\section{خلاصه}
\label{sec:math-book-summary}

در این فصل، یک فصل کامل از یک کتاب ریاضی
را دیدیم که شامل:

\begin{itemize}
    \item \textbf{تعاریف}: دنباله، همگرایی.
    \item \textbf{قضایا}: یکتایی حد، فشردگی،
    بولتزانو-وایرشتراس.
    \item \textbf{اثبات‌ها}: با استدلال ریاضی.
    \item \textbf{مثال‌ها}: دنباله‌ی هارمونیک.
    \item \textbf{نکته‌ها}: کاربرد قضیه‌ی
    فشردگی.
    \item \textbf{نمودارها}: همگرایی دنباله.
    \item \textbf{تمرین‌ها}: با راه‌حل.
    \item \textbf{نمایه}: با \cmd{index}.
    \item \textbf{ارجاعات}: با \cmd{cref}.
\end{itemize}

\mnote{این فصل، الگویی برای نوشتن فصل‌های کتاب ریاضی شماست. می‌توانید ساختار آن را برای موضوعات دیگر تطبیق دهید.}

در فصل بعدی (\cref{chap:programming-book})،
یک فصل نمونه از یک کتاب برنامه‌نویسی را
می‌بینیم که بر کد و الگوریتم تمرکز دارد.